\ifx\allfiles\undefined
\documentclass[12pt, a4paper, oneside, UTF8]{ctexbook}
\def\path{../config}
\input{\path/package.tex}

% 定理环境设置
\input{\path/theorem1_zh.tex}

\input{\path/custom.tex}

% 修改参数改变封面样式，0 默认原始封面、内置其他1、2、3种封面样式
\def\myIndex{3}


\ifnum\myIndex>0
    \input{\path/cover_package_\myIndex} 
\fi

\def\myTitle{高等数学笔记}
\def\myAuthor{M.K.}
\def\myDateCover{\today}
\def\myDateForeword{\today}
\def\myForeword{前言}
\def\myForewordText{
    
    这是一个基于\LaTeX{}模板的高数笔记．
    内容参考了包括但不限于以下资料：
    \begin{itemize}
        \item 《高等数学》（同济大学出版社）第七版
        \item 《吉米多维奇高等数学习题精选精讲》（山东科学技术出版社）第二版
        \item 《高等数学作业集》（哈尔滨工业大学出版社）第一版
        \item \href{https://space.bilibili.com/1035929235}{B站UP主：一高数}的高数课程
        \item \href{https://space.bilibili.com/42428180}{B站UP主：考研竞赛数学}的试题讲解与高数课程
        \item \href{https://space.bilibili.com/391930545}{B站UP主：Maki的完美算数教室}的高数内容

    \end{itemize}

    封面来源于\href{https://www.pixiv.net/users/3952}{おにねこ}的作品\href{https://www.pixiv.net/artworks/136551599}{雨雫とプレアデス}
    如有侵权请联系我删除．
}
\def\mySubheading{副标题}


\begin{document}
% \input{../config/cover}
\else
\fi
\chapter{二重极限}
\section{二重极限不存在的证明}
\begin{theorem}[二重极限不存在的常见证法]
一般的，证明二重极限不存在的方法有以下几种：
    \begin{enumerate}
        \item \textbf{路径法}：通过选择不同的路径（如$x$轴、$y$轴、$y=kx$等）来计算极限，如果不同路径得到的极限值不同或不存在，则二重极限不存在．
        \item \textbf{$\bm{y=kx^p}$路径法}：换元将式子凑成齐次，若结果依赖于$k$，则二重极限不存在．
    \item \textbf{极坐标法}：将$(x,y)$转换为极坐标$(\rho,\theta)$，如果极限值依赖于$\theta$，则二重极限不存在．
    \end{enumerate}
\end{theorem}


\begin{example}
    证明下列极限不存在：
    \begin{tasks}[label=(\arabic*), item-indent=1em](2)
        \task $\lim\limits_{(x,y) \to (0,0)} \dfrac{x+y}{xy}$
        \task $\lim\limits_{(x,y) \to (0,0)} \dfrac{xy}{x^2 + y^2}$
        \task $\lim\limits_{(x,y) \to (0,0)} \dfrac{xy^2}{x^2 + y^4}$
        \task $\lim\limits_{(x,y) \to (0,0)} \dfrac{x+y}{x-y}$
    \end{tasks}
\end{example}


\begin{proof}
    \begin{enumerate}[label=(\arabic*)]
        \item 单路径法：沿$y=x$趋近，$\lim\limits_{\substack{x \to 0 \\ y = x \to 0}} \dfrac{2x}{x^2} = \lim\limits_{x \to 0} \dfrac{2}{x} = \infty$；
        \item 双路径法：$\lim\limits_{ \substack{x \to 0 \\ y = x \to 0} } \dfrac{x^2}{x^2 + x^2} = \dfrac{1}{2}$；$\lim\limits_{ \substack{x \to 0 \\ y = -x \to 0} } \dfrac{-x^2}{x^2 + x^2} = -\dfrac{1}{2}$．不同路径得到的极限值不同，因此极限不存在．\\
            $y=kx$路径法：沿$y=kx$趋近，$\lim\limits_{x \to 0} \dfrac{x \cdot kx}{x^2 + (kx)^2} = \lim\limits_{x \to 0} \dfrac{k x^2}{x^2 (1 + k^2)} = \dfrac{k}{1 + k^2}$．$\lim\limits_{x \to 0} \dfrac{k}{1 + k^2}$依赖于$k$，因此极限不存在．\\
            极坐标法：将$(x,y)$转换为极坐标$(\rho,\theta)$，$\dfrac{xy}{x^2 + y^2} = \dfrac{\rho \cos\theta \cdot \rho \sin\theta}{\rho^2} = \cos\theta \sin\theta = \dfrac{1}{2} \sin 2\theta$．$\lim\limits_{\rho \to 0} \dfrac{1}{2} \sin 2\theta$依赖于$\theta$，因此极限不存在．
        \item $y=kx^p$路径法：令$x=ky^2$，则$\lim\limits_{y \to 0} \dfrac{k y^4}{k^2 y^4 + y^4} = \dfrac{k}{k^2 + 1}$．$\lim\limits_{y \to 0} \dfrac{k}{k^2 + 1}$依赖于$k$，因此极限不存在．
        \item 路径法：沿$y=x+x^2$趋近，$\lim\limits_{\substack{x \to 0 \\ y = x + x^2 \to 0}} \dfrac{x+x + x^2}{x-(x + x^2)} = \lim\limits_{x \to 0} \dfrac{2x+x^2}{-x^2} = \infty$
    \end{enumerate}
\end{proof}


\begin{myRemark}
    形如$\lim\limits_{(x,y) \to (0,0)} \dfrac{x^p y^q}{x^m+y^n}(m,n \in \mathbb{N^*}, p,q > 0)$的极限，有以下结论：
    \begin{enumerate}
        \item $m,n$中有奇数时，极限不存在；
        \item $m,n$中全为偶数时，\begin{itemize}
            \item 当$\dfrac{p}{m} + \dfrac{q}{n} > 1$时，极限为0；
            \item 当$\dfrac{p}{m} + \dfrac{q}{n} \le 1$时，极限不存在，可选取$y=kx^{\frac{m-p}{q}}$路径证明．
        \end{itemize}
    \end{enumerate}
\end{myRemark}


\section{二重极限的计算}
\begin{theorem}[二重极限的常见计算方法]
一般的，计算二重极限的方法有以下几种：
    \begin{enumerate}
        \item \textbf{连续性}：直接代入
        \item \textbf{等价无穷小}
        \item \textbf{夹逼定理}
        \item \textbf{$\bm{0\cdot}$ 有界量}
        \item \textbf{极坐标法}
    \end{enumerate}

\end{theorem}


\begin{example}
    计算下列极限：
    \begin{tasks}[label=(\arabic*), item-indent=1em](3)
        \task $\lim\limits_{(x,y) \to (0,0)} \dfrac{xy}{\sqrt{x^2 + y^2}}$
        \task $\lim\limits_{(x,y) \to (0,0)} \dfrac{\ln(1+x^2y^4)}{x^2 + y^2}$
        \task $\lim\limits_{(x,y) \to (0,0)} \dfrac{x^2+y^2}{|x|+|y|}$
    \end{tasks}
\end{example}


\begin{solution}
    \begin{enumerate}[label=(\arabic*)]
        \item 由于 $0 < \left| \dfrac{xy}{\sqrt{x^2 + y^2}} \right| \le |x| \to 0 \, (x \to 0)$，因此 $\lim\limits_{(x,y) \to (0,0)} \dfrac{xy}{\sqrt{x^2 + y^2}} = 0$．

        \item 
        \begin{align*}
            \lim\limits_{(x,y) \to (0,0)} \dfrac{\ln(1+x^2y^4)}{x^2 + y^2} &= \lim\limits_{(x,y) \to (0,0)} \dfrac{x^2y^4}{x^2 + y^2} \\
            &= \lim\limits_{(x,y) \to (0,0)} \dfrac{x^2}{x^2 + y^2} \cdot y^4
        \end{align*}

        由于 $\dfrac{x^2}{x^2 + y^2} \le 1$，因此 $\dfrac{x^2}{x^2 + y^2}$ 是有界量．又由于 $y^4 \to 0 \, (y \to 0)$，因此 
        \[\lim\limits_{(x,y) \to (0,0)} \dfrac{\ln(1+x^2y^4)}{x^2 + y^2} = 0\]

        \item 法一：原式可变形为 $|x| \dfrac{|x|}{|x|+|y|} + |y| \dfrac{|y|}{|x|+|y|}$，由于 $0 \le \dfrac{|x|}{|x|+|y|} \le 1$，$0 \le \dfrac{|y|}{|x|+|y|} \le 1$，因此 $\dfrac{|x|}{|x|+|y|}$ 和 $\dfrac{|y|}{|x|+|y|}$ 是有界量．又由于 $|x| \to 0 \, (x \to 0)$，$|y| \to 0 \, (y \to 0)$，因此 $\lim\limits_{(x,y) \to (0,0)} \dfrac{x^2+y^2}{|x|+|y|} = 0$．
        法二：令 $x = \rho \cos\theta$，$y = \rho \sin\theta$，则原式可变形为 $\dfrac{\rho^2}{\rho(|\cos\theta| + |\sin\theta|)} = \dfrac{\rho}{|\cos\theta| + |\sin\theta|}$．由于 $\dfrac{1}{|\cos\theta| + |\sin\theta|}$ 是有界量，$\rho \to 0 \, (\rho \to 0)$，因此 $\lim\limits_{(x,y) \to (0,0)} \dfrac{x^2+y^2}{|x|+|y|} = 0$．
    \end{enumerate}
\end{solution}


\begin{example}
    判断函数$f(x,y) = \begin{cases}
        \dfrac{2xy^2}{x^2+y^2}-\dfrac{2x^3y^2}{(x^2+y^2)^2}, & (x,y) \ne (0,0) \\ 0, & (x,y) = (0,0)
    \end{cases}$在点$(0,0)$处是否连续．
\end{example}

\
\begin{solution}
    \[
        \lim_{(x,y) \to (0,0)} f(x,y) = \lim_{(x,y) \to (0,0)} 2x\dfrac{y^2}{x^2+y^2} - x\dfrac{2x^2y^2}{(x^2+y^2)^2}
    \]
    又
    \[
        0 \le \dfrac{y^2} {x^2+y^2} \le 1, \quad 0 \le \dfrac{2x^2y^2}{(x^2+y^2)^2} \le 1
    \]
    因此，$\dfrac{y^2}{x^2+y^2}$和$\dfrac{2x^2y^2}{(x^2+y^2)^2}$是有界量．又由于$x \to 0 (x \to 0)$，因此$\lim_{(x,y) \to (0,0)} f(x,y) = 0$．又$f(0,0) = 0$，因此函数$f(x,y)$在点$(0,0)$处连续．
\end{solution}


\subsection{二重极限与二次极限的辨析}
\begin{example}
    求极限$\lim\limits_{x \to 0} \lim\limits_{y \to 0} x\sin\dfrac{1}{y} + y\sin\dfrac{1}{x}$
\end{example}

\begin{solution}
    由于$\lim\limits_{y \to 0} x\sin\dfrac{1}{y}$不存在，因此整个二次极限不存在．
\end{solution}
\ifx\allfiles\undefined
\end{document}
\fi